% This file is a chapter fragment.  It is included by version/detailed.tex, and
% it is inputted first, before every other chapter.  It is prose only: it
% carries no formalisation tags -- no \texttt{lean}, no \texttt{leanok}, no
% \texttt{uses} -- because it is not a formalisation target but a statement of
% the conventions the formalisation already follows.

\chapter{Conventions on coefficients}\label{sec:conventions-coefficients}

This chapter fixes, once, the coefficient framework of the project, so that a
reader of the blueprint can tell which of the several accounts collected here
is the one the Lean development implements.  Nothing in this chapter is
proved; each item below is a convention read off the Lean sources, and the
files that fix it are named.

\section{Sheaves on the small \'etale site}

A sheaf in the detailed chapters is a sheaf on the \emph{small} \'etale site of
a scheme $X$, with values in a coefficient category $A$.  In Lean this is
\texttt{Sheaf X.smallEtaleTopology A}; the small \'etale topology is mathlib's
\texttt{AlgebraicGeometry.Scheme.smallEtaleTopology}, and its objects are the
schemes \'etale over $X$.  Sheaves are formed on this site and on no other:
the pro-\'etale site and the fpqc site do not occur in the Lean development at
all.  Where an \'etale cover of $X$ is needed --- to say that a sheaf is
locally constant, for instance --- it is taken to be a cover in the big
\'etale precoverage \texttt{AlgebraicGeometry.Scheme.etalePrecoverage}, that
is, a jointly surjective family of \'etale maps $U_i\to X$; this names the
covering families, and no sheaf is ever formed on the big site.

The coefficient category $A$ is not a fixed category in Lean but a parameter,
required to be Grothendieck abelian and, following the convention of the
\texttt{SixFunctors} folder, to have $\operatorname{Hom}$-types in the universe
$u+1$, one above the universe $u$ of the base schemes.  Predicates involving
finiteness require $A$ to be moreover concrete, finiteness referring to the
underlying type.  This convention is stated in the module docstrings of
\begin{itemize}
  \item \texttt{WeilConjectures/SixFunctors/UpperLowerStar/Etale.lean}, which
        constructs the pushforward $g_*$ and the pullback $g^*$ between
        categories of sheaves on small \'etale sites, together with the
        adjunction $g^*\dashv g_*$, the composition isomorphisms and the
        exactness statements; and
  \item \texttt{WeilConjectures/ConstructibleSheaf/Defs.lean}, which defines
        the locally constant, finite, finite locally constant and
        constructible sheaves on the small \'etale site.
\end{itemize}
Both files record that the Grothendieck abelian generality is what supplies
sheafification and the left Kan extensions along base change used to define
$g^*$, and that some of those existence statements are, with current mathlib,
sorried instances rather than proved ones --- the small \'etale site is not
known to be essentially small, only the small affine \'etale site is.  This is
a limitation of the ambient library, not a mathematical hypothesis of the
theory.

\section{The base case is finite torsion coefficients}

The coefficient ring of the $\ell$-adic formalism is denoted $\Lambda$, and the
case the detailed chapters are written for is the \emph{finite torsion} one:
$\Lambda=\Lambda_n=\mathbb Z/\ell^n\mathbb Z$, so that $A$ is \texttt{ModuleCat}
$\Lambda$, the category of $\Lambda$-modules.  This is the instantiation named in the
docstring of \texttt{WeilConjectures/ConstructibleSheaf/Defs.lean}, both for
the coefficient category and for the constructibility predicate, and it is the
level at which every statement in the following chapters is to be read.  The
chapter on proper base change fixes a complete discrete valuation ring
$\Lambda$ with maximal ideal $\mathfrak m$ and writes
$\Lambda_n=\Lambda/\mathfrak m^n$; the finite levels $\Lambda_n$ are the finite
torsion rings above, and it is the finite-level statements that correspond to
the Lean code.

Where a chapter needs $\ell$ to be invertible, that is an assumption on the
schemes considered, not a change of coefficients.

\section{Adic coefficients are a later layer}

The passage from finite torsion coefficients to $\mathbb Z_\ell$, $\mathbb
Q_\ell$ and $\overline{\mathbb Q}_\ell$ --- derived completion, inversion of
$\ell$, and the colimit over finite extensions of $\mathbb Q_\ell$ --- is a
layer built on top of the finite torsion theory.  It is assumed nowhere in the
chapters below, and it is not implemented in Lean: the pro-\'etale site does
not occur in the development, and the coefficients there are always modules
over a ring, never a completed or adic object.  In the blueprint it lives
in the proper base change sections S8 (\'etale-to-pro-\'etale finite-level
transport), S9 (derived completion and inverse limits) and S10 (complete
pro-\'etale proper base change) of \Cref{chap:proper-base-change}, and in the
sketch articles 40--42 (\texttt{40\_finite\_torsion\_proetale.tex},
\texttt{41\_adic\_rational\_coefficients.tex},
\texttt{42\_qbarl\_triangulated\_colimit.tex}).  A statement in one of those
places is therefore not a statement about the objects the earlier chapters
formalise, and should not be cited as one.

\section{The \texttt{math/} folder is non-normative}

The folder \texttt{blueprint/src/version/math/} contains a separate account of
the theory --- the ``Route~B'' proof spine, which starts from the presentable
category of solid pro-\'etale complexes and builds the six operations there,
taking the large theorems as black boxes.  It is kept as background reading.
It is \emph{not} the convention of this project: it uses a different site and
different coefficients from the ones fixed above, and the Lean development does
not implement it.  Nothing in the detailed chapters depends on it.
